\documentclass[11pt]{article}
\usepackage[margin=1in]{geometry}
\usepackage{amsmath,amssymb,amsthm,booktabs,enumitem}
\usepackage{xcolor}
\usepackage[colorlinks=true,linkcolor=blue!50!black,citecolor=blue!50!black]{hyperref}
\newtheorem{proposition}{Proposition}
\newtheorem{lemma}[proposition]{Lemma}
\newtheorem{theorem}[proposition]{Theorem}
\theoremstyle{remark}
\newtheorem{remark}[proposition]{Remark}
\newcommand{\OO}{\mathcal{O}}
\newcommand{\Hc}{\mathcal{H}}
\setlist{itemsep=2pt}

\title{A worked example for $PGL(2)$\\[1ex]
\large Draft section for the 2026 revision of D.~Joyner,\\ ``Simple local trace formulas for unramified $p$-adic groups'' (1997)}
\author{Notes prepared by Claude (Anthropic) for D.~Joyner}
\date{November 10, 2026}

\begin{document}
\maketitle

\begin{abstract}
For $G=PGL(2)$ over a $p$-adic field we compute, in closed form, the distribution $I(\chi,\gamma)$ of the main identity
\[
\int_{T_\gamma\backslash G(F)} f(g^{-1}\gamma g)\,dg=\int f^\vee(\chi)\,I(\chi,\gamma)\,d\mu(\chi)\qquad(f\in\Hc(G,K),\ \gamma\in G_{ell})
\]
for every type of regular elliptic $\gamma$. The computation uses the Bruhat--Tits tree on the geometric side and Macdonald's formula on the spectral side. Every formula was checked numerically for $q=3$, against a direct count of vertices on the tree and by numerical integration over the unitary circle. The example confirms the corrected normalizations in the referee report (M1). It shows that $I(\chi,\gamma)$ is a rational function of $\chi$ whose poles at the non-regular characters are cancelled by the Plancherel density. It also shows that the dependence on the truncation parameter is a ``vertex term'' that pairs to zero \emph{exactly} with $f^\vee$ once the truncation is large. That last point is the model for the Route~A repair of \S4.
\end{abstract}

\section{Setup and conventions}

Let $F$ be a $p$-adic field with ring of integers $\OO$, uniformizer $\pi$ and residue field of order $q$. Let $G=PGL(2)$, $K=PGL(2,\OO)$, and let $A$ be the diagonal torus (modulo the center). Then $X_*(A)=\mathbb{Z}$, generated by the fundamental coweight $\varpi^\vee$ with $\varpi^\vee(\pi)=a_1$, where
\[
a_n:=\mathrm{diag}(\pi^n,1)\qquad(n\in\mathbb{Z}).
\]
The simple coroot is $\alpha^\vee=2\varpi^\vee$, so $a_\alpha=\alpha^\vee(\pi)=a_2$ modulo the center. The group $G$ is adjoint, which is why $X_*(A)\neq\mathbb{Z}\alpha^\vee$ here.

\paragraph{Characters, Weyl group, chamber.}
An unramified character $\chi$ of $A(F)$ is determined by $z=\chi(a_1)\in\mathbb{C}^\times$; it is unitary iff $|z|=1$. Then $\chi(a_n)=z^n$ and $\chi(a_\alpha)=z^2=:y$. The Weyl group $W=\{1,s\}$ acts by $z\mapsto z^{-1}$. The chamber $A^-=\{a:|\alpha(a)|\le1\}$ of the paper is $\{a_n:n\ge0\}$, since $|\alpha(a_n)|=|\pi^n|$. We have $\delta_B(a_n)^{1/2}=q^{-n/2}$.

\paragraph{Measures.}
We take $\mathrm{vol}(K)=1$. The centralizer $T_\gamma$ of a regular elliptic $\gamma$ is compact, and we give it volume~$1$. Then
\[
\OO_\gamma(f):=\int_{T_\gamma\backslash G}f(g^{-1}\gamma g)\,dg=\int_G f(g^{-1}\gamma g)\,dg.
\]
The double cosets are $Ka_nK$, $n\ge0$, with $\mathrm{vol}(Ka_nK)=(q+1)q^{n-1}$ for $n\ge1$. Let $f_n:=\mathbf 1_{Ka_nK}$.

\section{The spherical side}

\paragraph{Macdonald's formula.}
Let $c_\alpha(\chi)=\dfrac{1-q^{-1}\chi(a_\alpha)}{1-\chi(a_\alpha)}$ be as in (2.8) of the paper (split case). For $n\ge0$,
\begin{equation}\label{eq:mac}
\Gamma_\chi(a_n)=Q^{-1}q^{-n/2}\Big(c_\alpha(\chi^{-1})\,z^{n}+c_\alpha(\chi)\,z^{-n}\Big),\qquad Q=1+q^{-1}.
\end{equation}
Two features of this formula differ from Lemma~2.1 as printed:
\begin{itemize}
\item $Q=\sum_{w}q^{-\ell(w)}$, not $\sum_w(IwI:I)$ (report M1(a)).
\item On the paper's chamber $A^-$, the coefficient of $(w\chi)(a)$ is $c(w\chi^{-1})$, not $c(w\chi)$. In other words, Lemma~2.1 as printed is Macdonald's formula on the \emph{opposite} chamber.
\end{itemize}
Checks: at $n=0$ the bracket equals $1+q^{-1}$, so $\Gamma_\chi(1)=1$. Formula \eqref{eq:mac} was also verified numerically to machine precision.

\paragraph{Satake transform and Plancherel measure.}
The Satake transform of $f_n$ is
\[f_n^\vee(\chi)=\mathrm{vol}(Ka_nK)\,\Gamma_\chi(a_n).\] On the unitary circle $|z|=1$ (with $d\theta/2\pi$, $z=e^{i\theta}$), the Plancherel density is
\begin{equation}\label{eq:planch}
d\mu(\chi)=\rho(z)\,\frac{d\theta}{2\pi},\qquad
\rho(z)=\frac{Q}{|W|}\,|c_\alpha(\chi)|^{-2}=\frac{Q}{2}\,\frac{|1-y|^2}{|1-q^{-1}y|^2},\qquad y=z^2.
\end{equation}
\begin{proposition}[Plancherel inversion; Corollary 2.3]\label{prop:inv}
For all $n,m\ge0$,\[\int_{|z|=1}f_n^\vee(\chi)\,\Gamma_\chi(a_m)\,d\mu(\chi)=\delta_{nm}=f_n(a_m).\]
\end{proposition}
This was checked numerically for $0\le n,m\le3$ (Table~\ref{tab:inv}). With the printed constant $Q=1+q$, the diagonal entries would be $q$ instead of $1$.

\section{The geometric side: counting on the tree}

Let $\mathcal{X}$ be the Bruhat--Tits tree of $G$ (vertices $G/K$, regular of degree $q+1$) and let $d$ be its graph metric. Since $f_n$ is bi-$K$-invariant, $g\mapsto f_n(g^{-1}\gamma g)$ is right $K$-invariant, and $g^{-1}\gamma g\in Ka_nK$ iff $d(gK,\gamma gK)=n$. Hence
\begin{equation}\label{eq:count}
\OO_\gamma(f_n)=\#\{v\in\mathcal{X}:\ d(v,\gamma v)=n\}.
\end{equation}
The set is finite because $\gamma$ is elliptic.

Up to conjugacy and the center, a regular elliptic $\gamma$ generates $E=F[\gamma]$, a quadratic extension. There are three types:
\begin{description}[leftmargin=1.2em]
\item[(U$_r$)] $E/F$ unramified, $\gamma\in\OO_E^\times$, of depth $r\ge0$. The fixed-point set is the ball $B_r(v_0)$ of radius $r$ about a vertex. Each vertex at distance $m\ge1$ from the ball has $d(v,\gamma v)=2m$.
\item[(R$_r$)] $E/F$ ramified, $\gamma\in\OO_E^\times$, of depth $r\ge0$. The fixed-point set is the set of vertices within distance $r$ of an edge $e$. Each vertex at distance $m\ge1$ from it has $d(v,\gamma v)=2m$.
\item[(R$'$)] $E/F$ ramified and $v_E(\gamma)$ odd (e.g.\ $\gamma=\pi_E$). Then $\gamma$ reverses an edge $e$, and $d(v,\gamma v)=2\,d(v,e)+1$.
\end{description}
Counting vertices in spheres gives:
\begin{proposition}\label{prop:geom}
Let $m\ge1$.
\begin{center}
\begin{tabular}{lccc}
\toprule
type & $\OO_\gamma(f_0)$ & $\OO_\gamma(f_{2m})$ & $\OO_\gamma(f_{2m+1})$ \\
\midrule
(U$_r$) & $|B_r|=1+\dfrac{(q+1)(q^r-1)}{q-1}$ & $(q+1)q^{r+m-1}$ & $0$\\[2ex]
(R$_r$) & $\dfrac{2(q^{r+1}-1)}{q-1}$ & $2q^{r+m}$ & $0$\\[2ex]
(R$'$) & $0$ & $0$ & $2q^{m}$ \quad (also $\OO_\gamma(f_1)=2$)\\
\bottomrule
\end{tabular}
\end{center}
\end{proposition}
These counts were verified for $q=3$ by enumerating all $4373$ vertices of the ball of radius $7$ in the tree of $PGL(2,\mathbb{Q}_3)$, with $\gamma$ equal to
\[
\begin{pmatrix}0&-1\\1&0\end{pmatrix}\ (\mathrm{U}_0),\quad \begin{pmatrix}1&-3\\3&1\end{pmatrix}\ (\mathrm{U}_1),\quad \begin{pmatrix}1&3\\1&1\end{pmatrix}\ (\mathrm{R}_0),\quad \begin{pmatrix}0&3\\1&0\end{pmatrix}\ (\mathrm{R}').
\]
See Table~\ref{tab:data}.

\section{The spectral side: $I(\chi,\gamma)$ in closed form}

Write $f=\sum_n f(a_n)f_n$ (a finite sum), so $\OO_\gamma(f)=\sum_n f(a_n)\OO_\gamma(f_n)$. Insert Proposition~\ref{prop:inv} for $f(a_n)$. For any $N\ge\deg f$ (the largest $n$ with $f(a_n)\neq0$) this gives the \emph{exact} identity
\begin{equation}\label{eq:trunc}
\OO_\gamma(f)=\int f^\vee(\chi)\,I^N(\chi,\gamma)\,d\mu(\chi),\qquad I^N(\chi,\gamma):=\sum_{n=0}^{N}\OO_\gamma(f_n)\,\Gamma_\chi(a_n).
\end{equation}
This is the paper's truncated distribution $I^T$ of \S4.2, with $T\leftrightarrow N$. Now let $N\to\infty$. Write $y=z^2$ and $\tilde\theta(\chi)=1-\chi(a_\alpha)=1-y$.

\begin{theorem}\label{thm:main}
For regular unitary $\chi$ (i.e.\ $z\neq\pm1$), the limit $I(\chi,\gamma)=\lim_N I^N(\chi,\gamma)$ exists as an Abel limit. It is the rational function
\begin{align*}
\text{(U}_r\text{):}\quad I(\chi,\gamma)&=|B_r|\;+\;q^{r}\sum_{w\in W}c_\alpha\big((w\chi)^{-1}\big)\,\frac{(w\chi)(a_\alpha)}{\tilde\theta(w\chi)}
\;=\;|B_r|\;+\;q^{r}\,\frac{\cos\theta-q^{-1}}{1-\cos\theta}\qquad(y=e^{i\theta}),\\[1ex]
\text{(R}_r\text{):}\quad I(\chi,\gamma)&=\frac{2(q^{r+1}-1)}{q-1}\;+\;\frac{2q^{r+1}}{q+1}\cdot\frac{\cos\theta-q^{-1}}{1-\cos\theta},\\[1ex]
\text{(R}'\text{):}\quad I(\chi,\gamma)&=\frac{q^{1/2}(1-q^{-1})}{q+1}\cdot\frac{\cos\psi}{\sin^2\psi}\qquad(z=e^{i\psi}).
\end{align*}
Moreover, $\OO_\gamma(f)=\int f^\vee(\chi)I(\chi,\gamma)\,d\mu(\chi)$ for every $f\in\Hc(G,K)$, and the integral converges absolutely.
\end{theorem}

\begin{proof}[Sketch]
Consider type (U$_r$); the other types are identical. By Proposition~\ref{prop:geom} and \eqref{eq:mac}, the coefficient $\OO_\gamma(f_{2m})\,Q^{-1}q^{-m}$ equals $q^r$ for every $m\ge1$. So
\[
I^N=|B_r|+q^r\sum_{m=1}^{N/2}\big(c_\alpha(\chi^{-1})y^m+c_\alpha(\chi)y^{-m}\big).
\]
Summing the geometric series gives the closed form, plus the remainder
\[
I^N-I=-q^r\Big(c_\alpha(\chi^{-1})\frac{y^{M+1}}{1-y}+c_\alpha(\chi)\frac{y^{-M-1}}{1-y^{-1}}\Big),\qquad M=\lfloor N/2\rfloor.
\]
Multiply by the Plancherel density, using $\rho=\frac Q2\,(c_\alpha(\chi)c_\alpha(\chi^{-1}))^{-1}$ on $|z|=1$:
\[
(I^N-I)\,\rho=-\frac{Q}{2}\,q^{r}\Big(\frac{y^{M+1}}{1-q^{-1}y}+\frac{y^{-M-1}}{1-q^{-1}y^{-1}}\Big).
\]
The first term has only Fourier modes $y^k$ with $k\ge M+1$, the second only $k\le-M-1$. Since $f^\vee$ is a Laurent polynomial in $y^{\pm1/2}=z^{\pm1}$ of degree $\le\deg f$, we get $\int f^\vee(I^N-I)\,d\mu=0$ \emph{exactly} once $2(M+1)>\deg f$. Combine this with \eqref{eq:trunc}. Absolute convergence holds because $I\rho$ is smooth; see Remark~\ref{rem:poles}.
\end{proof}

\begin{remark}[Poles and the Plancherel density]\label{rem:poles}
$I(\chi,\gamma)$ has a \emph{double} pole at the non-regular characters $z=\pm1$ ($w\chi=\chi$), and no other poles on the unitary circle. The Plancherel density $\rho$ has a \emph{double zero} there. Their product is smooth, and in fact
\[
I(\chi,\gamma)\,\rho(z)=\frac{Q}{2}\cdot\frac{|B_r|\,(2-2\cos\theta)+q^r\,(2\cos\theta-2q^{-1})}{1-2q^{-1}\cos\theta+q^{-2}}\qquad\text{(type U}_r).
\]
So in rank one, ``rational function'' (1997 introduction) is right, \emph{provided} it is always paired with $d\mu$. On its own, $I(\cdot,\gamma)$ is not locally integrable at $z=\pm1$ against $d\theta$, which may explain why the 1992 abstract said ``distribution.''

The cancellation is structural. For each $w$,
\[
\frac{c_\alpha((w\chi)^{-1})\,\rho}{\tilde\theta(w\chi)}=\frac{Q}{2}\cdot\frac{1}{1-q^{-1}(w\chi)(a_\alpha)},
\]
because $\rho\propto|c|^{-2}$ supplies the factor $1-(w\chi)(a_\alpha)$ that cancels $\tilde\theta$. In higher rank the same mechanism cancels the factors of $\tilde\theta$ that vanish on walls. Whether cone denominators that vanish \emph{off} the walls also cancel, after summing over $w,w'$, is exactly the $(G,A)^\sim$-family statement of \S2.5. That is the one genuinely new ingredient Route~A needs in rank $\ge2$.
\end{remark}

\begin{remark}[What the example says about the truncation parameter]\label{rem:T}
Equation~\eqref{eq:trunc} shows that $f^\vee\mapsto\int f^\vee I^N\,d\mu$ is independent of $N$ once $N\ge\deg f$. That is automatic, and it is the first line of the paper's Theorem~4.5. What actually needs proof is \emph{regularity}: that $\lim_N I^N$ exists as a function and that $I\rho$ is integrable. The dependence on $N$ is a single ``vertex term'' $y^{M+1}\cdot(\text{holomorphic in }|y|<q)$, together with its $W$-conjugate, so it pairs to zero exactly, with no Riemann--Lebesgue argument needed. This suggests restating Proposition~4.2 and Theorem~4.5 as follows. Define $I(\chi,\phi)$ as the $T$-independent (vertex~$0$) part of the Brion decomposition of $I^T$. Then (i)~$I^T-I$ is a sum of vertex terms $\chi(T_\omega)R_\omega(\chi)$, and (ii)~each $R_\omega\rho$ is smooth with Fourier support in a translated cone. Together these give $\int f^\vee(I^T-I)\,d\mu=0$ for $T\ge T(f,\phi)$.
\end{remark}

\begin{remark}[The ``core'' term]
In type (U$_r$) the constant $|B_r|=\OO_\gamma(\mathbf 1_K)$ comes from the vertices fixed by $\gamma$. Those are exactly the points where $\gamma$ acts without displacement, so no Casselman asymptotic is available there. It is non-zero and depends on the depth $r$, while the second (cone) term has the shape predicted by Proposition~4.2: a sum over $W$ of $c$-functions over $\tilde\theta$. This is the concrete form of the gap noted in the log (item 3.69): any general argument must account for such a term, and the induction in Proposition~3.6 does not produce it. For type (R$'$) the core is empty, and $I$ is a pure cone term.
\end{remark}

\begin{remark}[Shape and sign]
$I(\chi,\gamma)$ depends on $\gamma$ only through its type and depth, as it must. In types (U$_r$) and (R$_r$), $I$ is positive on the regular unitary characters. Its minimum is at $\cos\theta=-1$ ($z=\pm i$), where for (U$_r$) it equals $|B_r|-q^r\frac{1+q^{-1}}{2}>0$, and $I\to+\infty$ as $z\to\pm1$. In type (R$'$), $I$ changes sign with $\cos\psi$. So $I(\cdot,\gamma)\,d\mu$ is a signed measure in general, not a positive one.
\end{remark}

\section{Numerical verification ($q=3$)}

All numbers below were produced by the script \texttt{pgl2\_check.py}, which accompanies this note. Orbital integrals are exact counts on the tree; spectral integrals use the midpoint rule with $4\times10^5$ nodes on $|z|=1$.

\begin{table}[h]
\centering
\begin{tabular}{c|cccc}
\toprule
$\int f_n^\vee\,\Gamma_\chi(a_m)\,d\mu$ & $m=0$ & $1$ & $2$ & $3$\\
\midrule
$n=0$ & 1.0000 & 0.0000 & 0.0000 & 0.0000\\
$n=1$ & 0.0000 & 1.0000 & 0.0000 & 0.0000\\
$n=2$ & 0.0000 & 0.0000 & 1.0000 & 0.0000\\
$n=3$ & 0.0000 & 0.0000 & 0.0000 & 1.0000\\
\bottomrule
\end{tabular}
\caption{Plancherel inversion with $Q=1+q^{-1}$ (Proposition~\ref{prop:inv}).}\label{tab:inv}
\end{table}

\begin{table}[h]
\centering
\begin{tabular}{c|cc|cc|cc|cc}
\toprule
& \multicolumn{2}{c|}{(U$_0$)} & \multicolumn{2}{c|}{(U$_1$)} & \multicolumn{2}{c|}{(R$_0$)} & \multicolumn{2}{c}{(R$'$)}\\
$n$ & tree & spectral & tree & spectral & tree & spectral & tree & spectral\\
\midrule
0 & 1 & 1.000000 & 5 & 5.000000 & 2 & 2.000000 & 0 & 0.000000\\
1 & 0 & 0.000000 & 0 & 0.000000 & 0 & 0.000000 & 2 & 2.000000\\
2 & 4 & 4.000000 & 12 & 12.000000 & 6 & 6.000000 & 0 & 0.000000\\
3 & 0 & 0.000000 & 0 & 0.000000 & 0 & 0.000000 & 6 & 6.000000\\
4 & 12 & 12.000000 & 36 & 36.000000 & 18 & 18.000000 & 0 & 0.000000\\
5 & 0 & 0.000000 & 0 & 0.000000 & 0 & 0.000000 & 18 & 18.000000\\
6 & 36 & 36.000000 & 108 & 108.000000 & 54 & 54.000000 & 0 & 0.000000\\
7 & 0 & 0.000000 & 0 & 0.000000 & 0 & 0.000000 & 54 & 54.000000\\
\bottomrule
\end{tabular}
\caption{$\OO_\gamma(f_n)$ counted on the tree versus $\int f_n^\vee I(\chi,\gamma)\,d\mu$ with $I$ from Theorem~\ref{thm:main}.}\label{tab:data}
\end{table}

We also checked the following, all to machine precision:
\begin{itemize}
\item the trigonometric forms in Theorem~\ref{thm:main};
\item the formula for $I\rho$ in Remark~\ref{rem:poles};
\item the chamber convention in \eqref{eq:mac}.
\end{itemize}
For the truncation (Remark~\ref{rem:T}) with $f=f_4$ and type (U$_0$): $\int f_4^\vee I^N\,d\mu$ equals $0$ for $N<4$ and exactly $12=\OO_\gamma(f_4)$ for all $N\ge4$.

\section{Consequences for the revision}
\begin{enumerate}
\item \textbf{Normalizations.} Fix Lemma~2.1 (constant $Q$, chamber $A^-$ with $c(w\chi^{-1})$) and Lemma~2.2 accordingly. Add these items to the log.
\item \textbf{Theorem 4.5.} State it with $I(\chi,\phi)$ defined as the $T$-independent part, and prove it by the exact-vanishing argument of Remark~\ref{rem:T}, not by ``independence of $T$ from Proposition~3.6.''
\item \textbf{Regularity.} In rank one, $I\,d\mu$ is smooth. In general this needs: (a)~the cancellation of wall factors (automatic, Remark~\ref{rem:poles}); and (b)~the $(G,A)^\sim$-family argument for cone denominators off the walls. Part~(b) is the main step of Route~A. A rank-two numerical test ($PGL(3)$, on the building) would be a good next check before writing it in general.
\item \textbf{Core term.} Proposition~3.6 must be restated as ``explicit cone terms $+$ a $T$-independent remainder.'' The example shows that the remainder is non-zero and geometric: $\OO_\gamma(\mathbf 1_K)$ in rank one.
\item \textbf{New section.} This example can go into the revision as a new section (or appendix), ``The case of $PGL(2)$,'' with the tables.
\end{enumerate}

\end{document}
